\chapter{The VNode model (what the renderer \emph{really} consumes)}

In Chapter~3 we treated TSX as ``just syntax'' that boils down to calling \texttt{h()} and \texttt{hFragment()}.
This chapter is where I make that concrete.

If you strip Relax down to its bones, the renderer is a pair of functions:

\begin{itemize}
  \item \texttt{mountDOM(vdom, host)}: allocate DOM for a VNode tree.
  \item \texttt{patchDOM(oldVdom, newVdom, host)}: reconcile two VNode trees (reuse DOM).
\end{itemize}

Both functions only work if their input has a stable, predictable data shape.
That data shape is the \textbf{VNode model} in \texttt{src/h.ts}.

I like to describe \texttt{src/h.ts} as ``the type-level constitution of the renderer''.
If you understand the invariants enforced here, every later chapter gets easier.

\section{What you will build in your head}
By the end of this chapter you should be able to do these three things confidently:

\begin{enumerate}
  \item Read a random VNode object in a debugger and know what it means.
  \item Predict what \texttt{h('div', \{\}, ['x'])} returns (including text node normalization).
  \item Explain why \texttt{areNodesEqual()} is written the way it is and how it drives patching.
\end{enumerate}

\section{The one invariant I defend everywhere}
Relax keeps the VNode model intentionally small, and then makes the runtime depend on one invariant:

\begin{quote}
  \noindent\textbf{Invariant:} after creation, every value inside \texttt{VNode.children} is already a VNode.
  No raw strings, no raw numbers, no \texttt{null}, no \texttt{undefined}.
\end{quote}

This is not a ``nice-to-have''.
It's the thing that keeps \texttt{mountDOM()} and \texttt{patchDOM()} from degenerating into dozens of slow type checks.

\section{Why we need VNodes in the first place}
The DOM is mutable state.
The renderer's job is to move a lot of mutable state toward a desired UI with minimal work.

So we need a representation of ``desired UI'' that is:

\begin{itemize}
  \item \textbf{Pure data:} can be created without touching the DOM.
  \item \textbf{Stable:} has explicit node kinds and predictable fields.
  \item \textbf{Diffable:} can be compared against a previous version.
\end{itemize}

In Relax, that's the VNode.

\section{Why VNodes exist}
A virtual DOM implementation needs a stable, serializable representation of UI. In Relax, that representation is the \textbf{VNode}. The renderer can then be split into three big phases:

\begin{enumerate}
  \item \textbf{Create}: build a VNode tree (usually by calling \texttt{h()} directly, or indirectly via TSX as described in Chapter~3).
  \item \textbf{Mount}: allocate real DOM nodes for that tree (\texttt{src/mount-dom.ts}).
  \item \textbf{Patch}: reconcile an ``old tree'' with a ``new tree'' (\texttt{src/patch-dom.ts}), reusing as much DOM as possible.
\end{enumerate}

The ``create'' step is deceptively important: the more regular the VNode tree is, the less branching the mount/patch code needs.

\subsection*{Diagram: from render to real DOM}
The following diagram is the "wiring diagram" for the rest of Part~II. Every later chapter is basically "zooming in" on one box:

\begin{center}
\begin{tikzpicture}[
  node distance=10mm,
  box/.style={draw,rounded corners,fill=RelaxLight,inner sep=5pt,align=center},
  arrow/.style={-Latex,thick}
]
  \node[box] (render) {\textbf{render()}\\(components / TSX)};
  \node[box, right=of render] (h) {\textbf{VNode factory}\\\texttt{h()}, \texttt{hFragment()}, \texttt{hSlot()}};
  \node[box, right=of h] (tree) {\textbf{VNode tree}\\(normalized children)};
  \node[box, below=of tree] (mount) {\textbf{Mount}\\\texttt{mountDOM}};
  \node[box, below=of h] (patch) {\textbf{Patch}\\\texttt{patchDOM}};
  \node[box, below=of render] (dom) {\textbf{Real DOM}\\nodes + listeners};

  \draw[arrow] (render) -- (h);
  \draw[arrow] (h) -- (tree);
  \draw[arrow] (tree) -- (mount);
  \draw[arrow] (mount) -- (dom);
  \draw[arrow] (tree) -- (patch);
  \draw[arrow] (patch) -- (dom);
\end{tikzpicture}
\end{center}

If you keep this picture in mind, a lot of micro-details (like \texttt{vdom.el} back-references or \texttt{areNodesEqual}) stop feeling ad-hoc: they're glue that keeps the arrows valid.

\section{The VNode types: \texttt{DOM\_TYPES}}
Relax defines node kinds as string constants in \texttt{DOM\_TYPES} (\texttt{src/h.ts}):

\begin{itemize}
  \item \texttt{TEXT}: a text node (\texttt{TextVNode}).
  \item \texttt{ELEMENT}: an HTML/SVG element (\texttt{ElementVNode}).
  \item \texttt{FRAGMENT}: a list of siblings without an extra wrapper element (\texttt{FragmentVNode}).
  \item \texttt{COMPONENT}: a component constructor + props/children (\texttt{ComponentVNode}).
  \item \texttt{SLOT}: a placeholder for ``external children'' in component trees (\texttt{SlotVNode}).
  \item \texttt{HRBR}: a special vnode wrapping an HRBR ``mount factory'' (\texttt{HrbrVNode}).
\end{itemize}

Two details are worth calling out because they influence the rest of the runtime:

\begin{enumerate}
  \item \textbf{Not every VNode has a DOM element.} \texttt{TEXT} uses a real \texttt{Text} node (\texttt{el?: Text}). \texttt{ELEMENT} uses a real \texttt{HTMLElement} (\texttt{el?: HTMLElement}). But \texttt{FRAGMENT} is special: it represents multiple siblings, so it can't point at a single ``root node''.
  \item \textbf{Keys are part of identity.} You'll see \texttt{props.key} used later by the reconciliation algorithm to decide whether a node is ``the same logical thing''.
\end{enumerate}

\section{Contract: \texttt{h(tag, props, children)}}
Relax's primary constructor is \texttt{h()} (\texttt{src/h.ts}). It's intentionally close to ``hyperscript'' functions found in other VDOM systems.

When I'm writing a renderer, I want the constructors to be boring.
``Boring'' here means: strongly asserted inputs, aggressively normalized children, minimal implicit behavior.
Relax follows that.

\subsection*{Code: the actual \texttt{h()} contract in \texttt{src/h.ts}}
The implementation is short enough that you should read it end-to-end at least once:

\begin{lstlisting}[language=JavaScript]
export function h(tag: string | unknown, props: Record<string, unknown> = {}, children: unknown[] = []): ElementVNode | ComponentVNode {
  const type = typeof tag === 'string' ? DOM_TYPES.ELEMENT : DOM_TYPES.COMPONENT

  assert(
    typeof props === 'object' && !Array.isArray(props),
    '[vdom] h() expects an object as props (2nd argument)'
  )
  assert(
    Array.isArray(children),
    `[vdom] h() expects an array of children (3rd argument), but got '${typeof children}'`
  )

  return {
    tag: tag as any,
    props: props as ElementVNodeProps,
    type,
    children: mapTextNodes(withoutNulls(children as any[])) as VNode[],
  } as any
}
\end{lstlisting}

Two things to notice:

\begin{itemize}
  \item \textbf{Fast fail:} \texttt{assert(...)} catches incorrect usage immediately.
  \item \textbf{Normalization happens right here:} \texttt{withoutNulls(...)} and \texttt{mapTextNodes(...)} enforce the children invariant.
\end{itemize}

\subsection*{Inputs}
\begin{itemize}
  \item \texttt{tag}: either a string (intrinsic element like \texttt{div}) or a component constructor/function.
  \item \texttt{props}: \emph{must be a plain object}. Relax checks this using \texttt{assert(...)} and fails fast with a helpful message.
  \item \texttt{children}: \emph{must be an array}. This is also asserted.
\end{itemize}

\subsection*{Output}
\begin{itemize}
  \item If \texttt{typeof tag === 'string'}, \texttt{h()} produces an \texttt{ELEMENT} vnode.
  \item Otherwise it produces a \texttt{COMPONENT} vnode.
\end{itemize}

Even though Relax supports TSX, TSX still ends up calling this layer (via the \texttt{relax-jsx} runtime), so understanding \texttt{h()} is foundational.

\section{Child normalization: \texttt{withoutNulls()} + \texttt{mapTextNodes()}}
Relax does two child normalization passes inside \texttt{h()} and \texttt{hFragment()}:

\begin{enumerate}
  \item \textbf{Drop nullish children}: \texttt{withoutNulls(children)} removes \texttt{null} and \texttt{undefined} (see \texttt{src/utils/arrays.ts}). This means ``conditional children'' can be written as \texttt{show ? node : null}.
  \item \textbf{Map primitives to \texttt{TEXT} nodes}: \texttt{mapTextNodes()} converts \texttt{string | number | boolean | bigint | symbol} into \texttt{TextVNode} via \texttt{hString()}.
\end{enumerate}

\subsection*{Why normalize eagerly?}
Normalization at creation time is a performance and correctness tactic:

\begin{itemize}
  \item \textbf{Performance:} mount/patch are on the hot path; you want them to pay as little "figure out what this is" cost as possible.
  \item \textbf{Correctness:} by forcing children into a single representation (VNodes), you avoid accidental differences between "first render" and "update" paths.
  \item \textbf{Debuggability:} tests can assert on VNode shapes without having to include JS coercion rules.
\end{itemize}

You can see this philosophy reflected in the tests: instead of testing dozens of mount/patch branches for primitives, the test suite locks down the normalization helpers and then relies on the invariant.

This is one of those design choices that pays rent everywhere:

\begin{itemize}
  \item Mount code can check \texttt{vNode.type} and never needs to special-case ``child is a string''.
  \item Patch code can treat text nodes uniformly; it doesn't have to worry about coercion rules.
\end{itemize}

\section{Fragments: \texttt{hFragment(children)} and \texttt{extractChildren()}}
Fragments represent ``multiple siblings'' and are heavily used by TSX fragment syntax (\texttt{<>...</>}).

\subsection*{Constructing fragments}
	exttt{hFragment(vNodes)} asserts the input is an array, filters nullish children, and maps primitives to text nodes (exactly like \texttt{h()}).

\subsection*{Code: fragment constructor}
This is the whole function (again: boring on purpose):

\begin{lstlisting}[language=JavaScript]
export function hFragment(vNodes: unknown[]): FragmentVNode {
  assert(Array.isArray(vNodes), '[vdom] hFragment() expects an array of vNodes')

  return {
    type: DOM_TYPES.FRAGMENT,
    children: mapTextNodes(withoutNulls(vNodes as any[])) as VNode[],
  }
}
\end{lstlisting}

\subsection*{Flattening fragments}
Later, other parts of the runtime sometimes need to treat ``a tree that may contain fragments'' as a flat list (for example, when calculating insertion positions).
	exttt{extractChildren(vdom)} provides that flattening:

\begin{itemize}
  \item If it finds a fragment child, it recursively collects that fragment's children.
  \item Otherwise it includes the child as-is.
\end{itemize}

The unit test \texttt{extract children from a tree with fragments} in \texttt{src/\_\_tests\_\_/h.test.ts} is your executable spec for this behavior.

\subsection*{Code: the spec is literally an object equality}
This is a good example of what I mean by ``tests as documentation'':

\begin{lstlisting}[language=JavaScript]
test('extract children from a tree with fragments', () => {
  const vNode = h('div', {}, [
    'A',
    hFragment([
      hFragment([hString('B')]),
      hString('C'),
      hFragment([hString('D')]),
    ]),
    'E',
  ])
  const children = extractChildren(vNode)

  expect(children).toEqual([
    { type: DOM_TYPES.TEXT, value: 'A' },
    { type: DOM_TYPES.TEXT, value: 'B' },
    { type: DOM_TYPES.TEXT, value: 'C' },
    { type: DOM_TYPES.TEXT, value: 'D' },
    { type: DOM_TYPES.TEXT, value: 'E' },
  ])
})
\end{lstlisting}

Notice what this test gives us: a guarantee that fragments behave like ``transparent containers'' for the purposes of child lists.

\section{Slots: \texttt{hSlot(children)} as a component-level placeholder}
Slots are how Relax expresses ``children passed into a component'' without forcing the component to take opaque \texttt{children} and do manual array plumbing.

At the VNode level, a slot is just:

\begin{itemize}
  \item \texttt{type: DOM\_TYPES.SLOT}
  \item optional \texttt{children}: default content to use when no external children are provided.
\end{itemize}

Then, to connect slot VNodes to component semantics, \texttt{src/h.ts} keeps a small piece of module-scoped state: \texttt{hSlotCalled}.

\subsection*{Why the module flag exists}
When rendering a component, Relax needs to know whether the component output contained at least one slot. Rather than scanning the render output tree after the fact, Relax uses this fast signal:

\begin{enumerate}
  \item \texttt{hSlot()} sets \texttt{hSlotCalled = true}.
  \item Component rendering checks \texttt{didCreateSlot()}.
  \item If a slot was created, the runtime calls the slot-filling algorithm (\texttt{src/slots.ts}) and then \texttt{resetDidCreateSlot()} to avoid leaking the flag into unrelated renders.
\end{enumerate}

This is one of Relax's ``tiny global state'' trade-offs: it keeps render fast and simple, but it also means slot creation is effectively a side-effect.

\subsection*{What the tests prove}
The best way to internalize slot semantics is to read \texttt{src/\_\_tests\_\_/component-slots.test.ts}. A few key cases:

\begin{itemize}
  \item \textbf{External content wins:} A component that renders \texttt{hSlot()} will receive the parent's children in that position.
  \item \textbf{Default content:} \texttt{hSlot([ ...default ])} is used when the parent provides no children.
  \item \textbf{Empty slot:} if there's no external content and no default content, the slot becomes empty (the test asserts \texttt{<div></div>}).
  \item \textbf{Conditional slots:} slot nodes can appear/disappear across re-renders, and patching handles it.
\end{itemize}

\section{HRBR interop: \texttt{hBlock(mount)} and VDOM identity}
HRBR (the compiled ``block'' runtime) needs to integrate into the same mounting/patching pipeline as regular VNodes. Relax does that by wrapping an HRBR mount factory in an \texttt{HRBR} vnode:

\begin{itemize}
  \item \texttt{hBlock(mount)} returns \texttt{\{ type: DOM\_TYPES.HRBR, mount \}}.
  \item Mounting creates a host element region and calls \texttt{mount(host)}.
  \item The returned instance exposes at least \texttt{destroy()}, and may expose \texttt{update()} / \texttt{dispose()}.
\end{itemize}

There's an important identity rule here, and the repo encodes it explicitly in \texttt{areNodesEqual()} (\texttt{src/nodes-equal.ts}):

\begin{quote}
  	extbf{HRBR vnode identity} is defined by mount factory identity: \texttt{old.mount === new.mount}.
\end{quote}

That single line is the difference between:

\begin{itemize}
  \item \textbf{Patch preserves instance:} stable mount function keeps the same HRBR instance.
  \item \textbf{Patch remounts:} new mount function destroys and mounts a fresh instance.
\end{itemize}

\subsection*{What the HRBR tests prove}
Read \texttt{src/\_\_tests\_\_/hrbr-integration.test.ts} as the contract:

\begin{itemize}
  \item \textbf{Mount and destroy:} an HRBR block returned from a VDOM component mounts into the DOM and gets destroyed on \texttt{destroyDOM}.
  \item \textbf{Event semantics:} HRBR event slots bind handler \texttt{this} to the host Relax component instance (the test asserts \texttt{calls[0].thisArg === vdom.component}).
  \item \textbf{Patch behavior:} stable vs. changing mount factory identity changes whether the HRBR instance is preserved.
\end{itemize}

This is a great example of the book's theme: treat tests as executable documentation.

\section{VNode identity and diffing: a preview}
This chapter is about creation, but VNode shape is designed with patching in mind. The function \texttt{areNodesEqual()} (\texttt{src/nodes-equal.ts}) shows you the identity rules Relax uses later:

\subsection*{Code: \texttt{areNodesEqual()} is tiny, but it's a design statement}
Read this function carefully.
Everything in patching depends on this decision about identity:

\begin{lstlisting}[language=JavaScript]
export function areNodesEqual(oldNode: any, newNode: any) {
  if (oldNode.type !== newNode.type) {
    return false
  }

  if (oldNode.type === DOM_TYPES.HRBR) {
    return oldNode.mount === newNode.mount
  }

  if (oldNode.type === DOM_TYPES.TEXT) {
    return true
  }

  if (oldNode.type === DOM_TYPES.ELEMENT) {
    return oldNode.tag === newNode.tag && oldNode.props?.key === newNode.props?.key
  }

  if (oldNode.type === DOM_TYPES.COMPONENT) {
    return oldNode.tag === newNode.tag && oldNode.props?.key === newNode.props?.key
  }

  return true
}
\end{lstlisting}

From a systems perspective, this is the policy wall between ``update in place'' and ``destroy + remount''.
If \texttt{areNodesEqual(old, next)} returns \texttt{false}, patching treats it as a different logical node and re-creates DOM / component instances.

\subsection*{Tests as spec: identity rules are enumerated}
The tests in \texttt{src/\_\_tests\_\_/node-equals.test.ts} are effectively a checklist:

\begin{itemize}
  \item Text nodes: always equal (value patches, identity doesn't change).
  \item Fragments: always equal.
  \item Elements: equal if tag matches, and key matches when present.
  \item Components: equal if constructor matches, and key matches when present.
\end{itemize}

This is why I keep the function small: it's easier to reason about and easier to lock down with tests.

\begin{itemize}
  \item text nodes: always equal (Relax patches text by updating value rather than replacing node type)
  \item fragments: always equal
  \item elements: equal if tag matches and \texttt{props.key} matches
  \item components: equal if constructor matches and \texttt{props.key} matches
  \item HRBR: equal if mount factory function is the same
\end{itemize}

The tests in \texttt{src/\_\_tests\_\_/node-equals.test.ts} enumerate these cases in small, sharp assertions.

\section{Mini-exercises}
\begin{enumerate}
  \item \textbf{Read like a maintainer:} In \texttt{src/\_\_tests\_\_/h.test.ts}, locate the test named \texttt{"h() maps 'number', 'boolean', 'bigint', and 'symbol' values to text vNodes"}. Explain (to yourself) why \texttt{Symbol(5)} becomes the string \texttt{"Symbol(5)"}.
  \item \textbf{Write a new invariant test:} Add a small test that asserts \texttt{h('div', {}, [null, undefined, 'x']).children.length === 1}. (This should already be true, but the test helps lock the behavior down.)
  \item \textbf{Identity thought experiment:} Why do you think Relax treats text nodes as ``always equal''? What does that allow \texttt{patchDOM()} to do as an optimization?
\end{enumerate}

\section{Online resources}
This chapter describes patterns that show up across many UI runtimes. These references help if you want to compare Relax's choices to other ecosystems:

\begin{itemize}
  \item Hyperscript and virtual node factories (conceptual background): \href{https://github.com/hyperhype/hyperscript}{https://github.com/hyperhype/hyperscript}
  \item DOM node types (Text vs Element) on MDN: \href{https://developer.mozilla.org/en-US/docs/Web/API/Node/nodeType}{https://developer.mozilla.org/en-US/docs/Web/API/Node/nodeType}
  \item React's reconciliation keys (why \texttt{key} exists): \href{https://react.dev/learn/rendering-lists}{https://react.dev/learn/rendering-lists}
  \item Web Components slots (a parallel mental model for ``default vs external'' content): \href{https://developer.mozilla.org/en-US/docs/Web/API/Web_components/Using_templates_and_slots}{\url{https://developer.mozilla.org/en-US/docs/Web/API/Web_components/Using_templates_and_slots}}
\end{itemize}
